\documentclass[10pt, aspectratio=169]{beamer}
\input{../../_en_preambulo_beamer}
% Comandos y macros estadísticas compartidas para las presentaciones Beamer en inglés (Metropolis)

% Conjuntos numéricos básicos
\newcommand{\R}{\mathbb{R}}
\newcommand{\N}{\mathbb{N}}
\newcommand{\Q}{\mathbb{Q}}
\newcommand{\Z}{\mathbb{Z}}
\newcommand{\set}[1]{\left\{ #1 \right\}}
\newcommand{\abs}[1]{\left|#1\right|}
\newcommand{\norm}[1]{\left\| #1 \right\|}

% Comandos estadísticos y probabilísticos del curso
\newcommand{\Var}{\operatorname{Var}}
\newcommand{\cov}{\operatorname{Cov}}
\newcommand{\corr}{\rho}
\newcommand{\comb}[2]{\begin{pmatrix} #1 \\ #2 \end{pmatrix}}
\newcommand{\s}{\sigma}
\newcommand{\particion}{\mathfrak{P}}
\newcommand{\card}[1]{\#\left( #1 \right)}
\newcommand{\seq}[3]{\set{#1}_{#2}^{#3}}
\newcommand{\E}{\mathbb{E}}
\newcommand{\Prob}{\mathbb{P}}

% Entornos pedagógicos y cajas en inglés académico natural (soportando tanto nombres en inglés como en español)
\newenvironment{discussion}{%
  \begin{alertblock}{Discussion Question}%
}{%
  \end{alertblock}%
}
\newenvironment{discusion}{%
  \begin{alertblock}{Discussion Question}%
}{%
  \end{alertblock}%
}

\newenvironment{reflection}{%
  \begin{alertblock}{Classroom Reflection}%
}{%
  \end{alertblock}%
}
\newenvironment{reflexion}{%
  \begin{alertblock}{Classroom Reflection}%
}{%
  \end{alertblock}%
}

\newenvironment{paradox}{%
  \begin{alertblock}{Exploratory Paradox}%
}{%
  \end{alertblock}%
}
\newenvironment{paradoja}{%
  \begin{alertblock}{Exploratory Paradox}%
}{%
  \end{alertblock}%
}

\newenvironment{motivation}{%
  \begin{exampleblock}{Motivation: Data Science in Practice}%
}{%
  \end{exampleblock}%
}
\newenvironment{motivacion}{%
  \begin{exampleblock}{Motivation: Data Science in Practice}%
}{%
  \end{exampleblock}%
}

\newenvironment{quickchallenge}{%
  \begin{block}{Quick Team Challenge}%
}{%
  \end{block}%
}
\newenvironment{retoclase}{%
  \begin{block}{Quick Team Challenge}%
}{%
  \end{block}%
}


\title[Introduction to Regression]{Section 09.02: Introduction to Linear Regression}
\subtitle{Stochastic Models, Classical Assumptions, and Applications}
\author[J. Castillo Colmenares]{Juliho Castillo Colmenares}
\institute{Tecnológico de Monterrey}
\date{\vspace{-1.5cm}}

\begin{document}

% Slide 1: Title
\begin{frame}[plain]
  \titlepage
\end{frame}

% Slide 2: Roadmap
\begin{frame}{Roadmap --- Chapter 09: Linear and Multiple Regressions}
  \scriptsize
  Where are we in the modular development of this chapter?
  \begin{itemize}\itemsep=0.02em
    \item \textbf{09.01} Correlation as the Premise of Regression
    \item \textbf{09.02} Introduction to Linear Regression \emph{(Today)}
    \item \textbf{09.03} The Mathematics of Regression: OLS Derivation and $R^2$
    \item \textbf{09.04} Linear Regression on Simulated Data
    \item \textbf{09.05} Optimal Coefficients, $t$/$F$ Tests, and RSE
    \item \textbf{09.06} Implementation in Python with \texttt{statsmodels}
    \item \textbf{09.07} Multiple Regression, the Normal Equation, and Ridge/Lasso
    \item \textbf{09.08} Model Validation and $k$-fold Cross-Validation
    \item \textbf{09.09} Residual Diagnostics and Multicollinearity (VIF)
    \item \textbf{09.10} Regression with \texttt{scikit-learn}
    \item \textbf{09.11} Categorical Variables and Dummy Variables
    \item \textbf{09.12} Nonlinear Transformations and Polynomial Regression
  \end{itemize}
  \pause
  \vspace{-0.05cm}
  \begin{block}{Session Objective}
    Present linear regression as a predictive stochastic model: distinguish it from deterministic models, formally differentiate it from correlation (Section 09.01), enumerate its five classical assumptions, and place simple and multiple regression within the landscape of data science applications.
  \end{block}
\end{frame}

% Slide 3: Motivation
\begin{frame}{From Association to Predictive Model}
  \footnotesize
  \begin{motivacion}
    Linear regression is a technique developed in the 19th century by Francis Galton and mathematically formalized by Karl Pearson and Ronald Fisher. Despite its simplicity, it serves as the foundation for more advanced techniques: polynomial regression, generalized linear models, and machine learning methods.
  \end{motivacion}

  \vspace{0.15cm}
  \pause
  \begin{itemize}\itemsep=0.1cm
    \item \textbf{Prediction:} estimating future values from historical data. \pause
    \item \textbf{Inference:} understanding and quantifying relationships between variables.
  \end{itemize}
\end{frame}

% Slide 4: Deterministic vs stochastic models
\begin{frame}{Deterministic vs. Stochastic Models}
  \footnotesize
  \begin{block}{Deterministic Model}
    Always produces the same output for the same inputs (e.g. $F=ma$): there is no error term.
  \end{block}
  \pause

  \vspace{0.1cm}
  \begin{alertblock}{Stochastic Model}
    Linear regression recognizes that real observations do not follow a linear relationship perfectly: it includes a random error component $\epsilon$, reflecting factors not captured by the model.
  \end{alertblock}
\end{frame}

% Slide 5: Model example
\begin{frame}{Example: House Price as a Linear Model}
  \footnotesize
  \begin{block}{Model Specification}
    If house price $P$ depends linearly on size $S$, amenities $A$, and transit access $T$:
    $$P = a_1 S + a_2 A + a_3 T + \epsilon$$
  \end{block}
  \pause

  \vspace{0.1cm}
  \begin{alertblock}{Components}
    $P$ is the output (predicted) variable; $S,A,T$ are the input (known) variables; $a_1,a_2,a_3$ are parameters that must be \textbf{estimated from historical data}; $\epsilon$ is the irreducible random error.
  \end{alertblock}
\end{frame}

% Slide 6: Correlation vs regression
\begin{frame}{Correlation vs. Regression: the Formal Distinction}
  \footnotesize
  \begin{block}{Correlation (Section 09.01)}
    Measures strength and direction; is \textbf{symmetric} ($r_{XY}=r_{YX}$) and dimensionless.
  \end{block}
  \pause

  \vspace{0.1cm}
  \begin{alertblock}{Regression}
    Models the functional relationship to \textbf{predict} values; is \textbf{not symmetric} (regressing $Y$ on $X$ gives different coefficients than regressing $X$ on $Y$) and its coefficients carry units.
  \end{alertblock}
\end{frame}

% Slide 7: Classical assumptions
\begin{frame}{The Five Classical Assumptions (Preview)}
  \footnotesize
  For a regression model's inferences to be valid:
  \begin{enumerate}\itemsep=0.06cm
    \item \textbf{Linearity} of the predictor-response relationship. \pause
    \item \textbf{Independence} of observations. \pause
    \item \textbf{Homoscedasticity}: constant error variance. \pause
    \item \textbf{Normality} of the errors. \pause
    \item \textbf{No multicollinearity} among predictors.
  \end{enumerate}
  \vspace{0.05cm}
  \begin{alertblock}{Note}
    The full formal audit of these five assumptions is developed in Section 09.09.
  \end{alertblock}
\end{frame}

% Slide 8: Types of regression
\begin{frame}{Simple vs. Multiple Regression}
  \footnotesize
  \begin{block}{Simple Linear Regression}
    A single predictor variable: $Y=\beta_0+\beta_1X+\epsilon$ (Sections 09.03-09.06).
  \end{block}
  \pause

  \vspace{0.1cm}
  \begin{alertblock}{Multiple Linear Regression}
    Two or more predictors: $Y=\beta_0+\beta_1X_1+\cdots+\beta_pX_p+\epsilon$ (Section 09.07). Both types share the same underlying principle: ordinary least squares.
  \end{alertblock}
\end{frame}

% Slide 9: Python Lab (Block 1)
\begin{frame}[fragile]{Python Lab: Deterministic vs. Stochastic}
  \scriptsize
  Zero residual variance in a deterministic model vs. nonzero variance in a stochastic one (\texttt{09.02\_introduction\_to\_regression.py}):
  {\lstset{basicstyle=\fontsize{3.6pt}{4.3pt}\selectfont\ttfamily,aboveskip=0.05em,belowskip=0.05em}
  \lstinputlisting[firstline=17, lastline=27]{../../code/09_regresiones/09.02_introduction_to_regression.py}}
\end{frame}

% Slide 10: Python Lab (Block 2)
\begin{frame}[fragile]{Python Lab: Simple vs. Multiple Regression}
  \scriptsize
  Comparing $R^2$ when going from one predictor to three, using the house-price example:
  {\lstset{basicstyle=\fontsize{3.2pt}{3.9pt}\selectfont\ttfamily,aboveskip=0.05em,belowskip=0.05em}
  \lstinputlisting[firstline=30, lastline=47]{../../code/09_regresiones/09.02_introduction_to_regression.py}}
\end{frame}

% Slide 11: Python Lab (Block 3)
\begin{frame}[fragile]{Python Lab: Assumptions Preview}
  \scriptsize
  Quick numeric check of zero-mean and homoscedasticity (full audit in 09.09):
  {\lstset{basicstyle=\fontsize{3.4pt}{4.1pt}\selectfont\ttfamily,aboveskip=0.05em,belowskip=0.05em}
  \lstinputlisting[firstline=50, lastline=64]{../../code/09_regresiones/09.02_introduction_to_regression.py}}
\end{frame}

% Slide 12: Terminal Output
\begin{frame}[fragile]{Python Lab: Terminal Output}
  \scriptsize
  Standard output produced when running the lab script:
  \vspace{0.02cm}
  \begin{block}{Standard Console Output}
    \fontsize{3.6pt}{4.3pt}\selectfont
    \begin{verbatim}
=== Block 1: Deterministic vs. Stochastic ===
Deterministic (F=ma):     residual variance = 0.000000
Stochastic (price~size):  residual variance = 65.9973

=== Block 2: Simple vs. Multiple Regression ===
Simple (price~size):                  R^2=0.6728
Multiple (price~size+amenities+transit): R^2=0.8893 (+0.2165)

=== Block 3: Assumptions Quick Preview ===
mean(residuals) = -0.000000 (zero-mean OK)
Var(first half)=2.35, Var(second half)=2.45 (homoscedasticity OK)
    \end{verbatim}
  \end{block}
\end{frame}

% Slide 13: Applications
\begin{frame}{Practical Applications of Linear Regression}
  \footnotesize
  \begin{itemize}\itemsep=0.1cm
    \item \textbf{Finance:} stock price prediction, credit risk assessment. \pause
    \item \textbf{Marketing:} modeling the impact of advertising on sales (the Section 09.01 case). \pause
    \item \textbf{Medicine:} relating drug dosages to clinical responses. \pause
    \item \textbf{Engineering and social sciences:} quality control, failure prediction, socioeconomic studies.
  \end{itemize}
\end{frame}

% Slide 14: Closing
\begin{frame}{Section Closing: Introduction to Linear Regression}
  \begin{reflexion}
    Linear regression formalizes the association measured by correlation into a predictive stochastic model: it explicitly recognizes that real observations include random error, and rests on five classical assumptions that guarantee the validity of its inferences.
  \end{reflexion}

  \vspace{0.2cm}
  \pause
  \begin{block}{Key Takeaways}
    \begin{itemize}\itemsep=0.08cm
      \item \textbf{Stochastic, not deterministic:} the $\epsilon$ term is central, not a minor detail.
      \item \textbf{Correlation $\neq$ regression:} regression predicts and its coefficients have units and asymmetric interpretation.
      \item \textbf{Simple vs. multiple:} same principle (least squares), different number of predictors.
    \end{itemize}
  \end{block}
\end{frame}

% Slide 15: Modular Perspective
\begin{frame}{Modular Perspective --- The Next Challenge}
  \scriptsize
  \begin{retoclase}
    With the model conceptually established, Section 09.03 answers the central mathematical question: how exactly are the parameters $\beta_0,\beta_1$ estimated from the data? We will rigorously derive the ordinary least squares (OLS) formulas and prove the fundamental decomposition $\text{SST}=\text{SSR}+\text{SSE}$ that gives rise to the coefficient $R^2$.
  \end{retoclase}

  \vspace{0.08cm}
  \pause
  \begin{alertblock}{Recommended Preparation}
    \begin{itemize}\itemsep=0.06cm
      \item Review basic differential calculus: minimizing a function via partial derivatives.
      \item Reflect on what "minimizing the sum of squared errors" means geometrically.
    \end{itemize}
  \end{alertblock}
\end{frame}

\end{document}
